\section{Proof of \cref{3.3.7}}\label{secB}

\begin{lemma}\label{B.1.1}
	Let $x\in\X$ be a degenerate simplex. Then there is a unique nondegenerate simplex $y\in X_{k}  $ for $k\leq n$, and a unique composite of codegeneracies $\alpha : [n] \to [k]$ in $\boldsymbol{\Delta} $ such that $\alpha^*(x)=y$, where $\alpha^*$ is the simplicial operator induced by $\alpha$ in the simplicial set $X_{\bullet} $.
\end{lemma}
\begin{proof}
	First we show existence. Let $x\in X_n$ be degenerate. By definition, there exists $y\in X_{n-1} $ and $i\leq n$ such that $s_i(y)=x$. If $y$ is nondegenerate, we are done. Otherwise, we continue inductively with the same argument. Note that there is an upper bound; in particular all $0$-simplices are nondegenerate since no degeneracies have domain $X_0$. Since $n$ is finite, induction gives us the required existence.

	Uniqueness of the map is obtained by noting any composite of degeneracies is a composite of codegeneracies in $\boldsymbol{\Delta} $, which is a surjective map, and as a subresult of \cref{2.3.4}, surjections are \emph{uniquely} determined by their degenerate decomposition. Uniqueness of the nondegenerate simplex is obtained as follows: let $\alpha^*(y)=\alpha^*(z)=x$. Since $\alpha^*$ is a composite of degeneracies, $\alpha$ is a surjection in $\boldsymbol{\Delta} $, which means it admits a right inverse by the axiom of choice. Denote this right inverse by $\sigma$. We prove by contradiction that $\sigma$ is order preserving. Note that $\sigma : [k]\to [n]$. If $k=0$, then $\sigma$ is trivially order preserving, so we may freely assume there exist distinct $a,b\in[k]$. In particular, we may let $a< b$ strictly, and suppose that $\sigma(a)\geq \sigma(b)$. Then since $\alpha$ is order preserving, $\alpha\sigma(a)\geq \alpha\sigma(b)$, but since $\alpha\sigma=1$, we have that $a\geq b$, a contradiction. Thus, $\sigma$ is order preserving, and a morphism in $\boldsymbol{\Delta} $. It induces a simplicial operator $\sigma^*$ in $X_{\bullet} $, which must satisfy
	\[
		z=(\alpha\sigma)^*(z)=\sigma^*\alpha^*(z)=\sigma^(x)=\sigma^*\alpha^*(y)=(\alpha\sigma)^*(y)=y
	.\]
	Thus, the nondegenerate simplex is unique.
\end{proof}

\begin{remark}\label{B.1.2}
	We would like to quotient only by faces, but the issue with quotienting only by $(d_i(x),p)\sim (x,d^i(p))$ is that faces of a nondegenerate simplex need not be nondegenerate (\cite{fried}; see the discussion following definition 3.8), and thus wouldn't be present in the new space, so the relation wouldn't be very well-defined. We instead pass to the relation $\sim _f$, which is defined as the relation generated by the following rule.

	Let $X^{nd} _n$ denote the set of \emph{nondegenerate} $n$-simplices of $X_{\bullet} $. Define the equivalence relation $\sim _f$ on $\coprod_{[n]\in\boldsymbol{\Delta} } X_n^{nd} \times \left| \Delta^n \right| $ by $(x,\mu_*(p))\sim _f (y,\beta_*(p))$ if
	\begin{itemize}
		\item $\mu$ is a composite of coface maps;
		\item the unique nondegenerate simplex of $\mu^*(x)$ given by \cref{B.1.1} is $y$;
		\item and $\beta$ is the unique composite of codegeneracies given by \cref{B.1.1} such that $\beta^*(y)=\mu^*(x)$.
	\end{itemize}
	Note that $\mu$ can be the identity by choosing it to be the empty composite, which is an empty composite of coface maps, giving
	\[
		(x,p)\sim_f (x,p),
	\]
	since $x$ is already nondegenerate and thus the corresponding composite of codegeneracies is also empty, hence another identity. Since the relation is reflexive, we simply demand symmetry and transitivity, giving us the equivalence relation by generating it on the above rule.
\end{remark}


\begin{proposition}[Prop. 3.3.7]\label{B.1.3}
	Let $X_{\bullet} $ be a simplicial set, let $\sim_f$ be as defined in \cref{B.1.2}, let $\sim $ be the relation given in \cref{3.2.2} (2), and let $X_n^{nd} $ denote the set of nondegenerate $n$-simplices of $X$. Then there is a homeomorphism
	\[
		\left| X_{\bullet}  \right| \coloneqq \left( \coprod_{n\geq 0} X_n\times \left| \Delta^n \right|  \right) /\tildesim \cong \left( \coprod_{n\geq 0} X_n^{nd} \times \left| \Delta^n \right|  \right) /\tildesim _f
	.\]
	In words, one may construct the geometric realization by deleting nondegenerate simplices, and quotienting by the relation of \cref{B.1.2}, which is extremely similar to quotienting purely by faces.
\end{proposition}
\begin{proof}
	We want to show
	\[
		\left( \coprod_{[n]\in\boldsymbol{\Delta} } X_n\times \left| \Delta^n \right|  \right) /\tildesim \cong \left( \coprod_{[n]\in\boldsymbol{\Delta} } X_n^{nd} \times \left| \Delta^n \right|  \right) /\tildesim_f ,
	\]
	Appealing to \cref{B.1.1}, we define the map
	\[\begin{tikzcd}[row sep = 0em, column sep = 2.5em]
		\varphi  : \displaystyle{\coprod_{[n]\in\boldsymbol{\Delta} } X_n\times \left| \Delta^n \right| }\ar[r]  & \displaystyle{\coprod_{[n]\in\boldsymbol{\Delta} } X_n^{nd} \times \left| \Delta^n \right|} \\
		(x,p)=(\alpha^*(y),p)\ar[mapsto, r]  & {(y,\alpha_*(p))},
	\end{tikzcd}\]
	where $\alpha^*$ is as defined in \cref{B.1.1}, and $\alpha_*=\left| \Delta^{\bullet}  \right| (\alpha): \left| \Delta^k \right|  \to \left| \Delta^n \right| $. We defer showing the map is continuous until the end of the proof. Our goal for now will be to show that if $x\sim y$ in the domain of $\varphi$, then $\varphi(x)\sim _f\varphi(y)$ in the codomain; postcomposing with the quotient morphism will then identify these elements, and thus the composite identifies all elements under $\sim $, hence we may invoke universality.

	First, we show it respects degeneracies. We can show a strict equality here, thus the relation follows by reflexivity. Let $x\in X_n$ and $p\in\left| \Delta^{n+1}  \right| $, such that $(s_i(x),p)\sim (x,s^i(p))$. By \cref{B.1.1}, let $y\in X_k$ be nondegenerate, and let $\alpha^*(y)=x$. By uniqueness of $\alpha$, we have that $s_i\alpha^*(y)$ is the unique composition of degeneracies mapping $y\mapsto x$. Note also that $s_i\alpha^*=X(\alpha s^i)$ by contravariance of $X_{\bullet} $, and thus the induced morphism on geometric simplices is $\left| \Delta^{\bullet}  \right| (\alpha s^i)=\alpha_*s^i$ by functoriality. We have
	\[
		\varphi(s_i(x),p)=\varphi(s^i\alpha^*(y),p)=(y,\alpha_*s^i(p)),\qquad\varphi(x,s^i(p))=\varphi(\alpha^*(y),s^i(p))=(y,\alpha_*s^i(p))
	.\]
	Thus, $\varphi$ respects $\sim $ for degeneracies.

	We now check faces. Recall that we must show that $\varphi(d_i(x),p)\sim_f\varphi(x,d^i(p))$. Keep notation as before, meaning $x=\alpha^*(y)$ is given by \cref{B.1.1}. We have
	\[
		\varphi(x,d^i(p))=\varphi(\alpha^*(y),d^i(p))=(y,\alpha_*d^i(p))
	.\]
	The other side is more difficult. Note that $(d_i(x),p)=(d_i\alpha^*(y),p)=((\alpha d^i)^*(y),p)$ by contravariance of $X_{\bullet} $. By \cref{2.3.4}, $\alpha d^i$ admits a unique factorization into a composite of codegeneracies (a epi component) and a composite of cofaces (a monic component). Write $\alpha d^i=\mu\eta$ for this monic-epi decomposition, where $\mu$ is the monic and $\eta$ is the epi. We thus have
	\[
		d_i(x)=d_i\alpha^*(y)=(\alpha d^i)^*(y)=(\mu\eta)^*(y)=\eta^*\mu^*(y)
	.\]
	Since $\mu^*(y)$ may be degenerate, we define $\mu^*(y)=\beta^*(z)$, where $z$ is the unique nondegenerate simplex and $\beta$ is the unique composite of degeneracies, as given in \cref{B.1.1}. By uniqueness, we have that $d_i(x)=\eta^*\beta^*(z)=(\beta\eta)^*(z)$ is the unique composite of degeneracies for $d_i(x)$ given by \cref{B.1.1}, so
	\[
		\varphi(d_i(x),p)=\varphi((\beta\eta)^*(z),p)=(z,\beta_*\eta_*(p))
	.\]
	Returning to $(x,d^i(p))$, we note that since $d^i\alpha=\mu\eta$, since $\beta^*(z)=\eta^*(y)$, and since $\mu$ is a composite of cofaces, we have
	\[
		\varphi(x,d^i(p))=(y,\alpha_*d^i(p))=(y,\mu_*\eta_*(p))\sim _f (z,\beta_*\eta_*(p))=\varphi(d_i(x),p)
	\]
	by definition of $\sim _f$ as given in \cref{B.1.2}.

	We have thus shown $(x,p)\sim (y,q)$ implies $\varphi(x,p)\sim_f \varphi(y,q)$. Since the quotient map $\pi_{nd} $ equalizes $\sim _f$, have that $\pi_{nd} \varphi$ descends along the quotient
	\[\begin{tikzcd}
	{\displaystyle{\coprod_{[n]\in\boldsymbol{\Delta}}X_n\times\left|\Delta^n\right|}} & {\displaystyle{\coprod_{[n]\in\boldsymbol{\Delta}}X_n^{nd}\times\left|\Delta^n\right|}} \\
	{\displaystyle{\left(\coprod_{[n]\in\boldsymbol{\Delta}}X_n\times\left|\Delta^n\right|\right)/\tildesim}} & {\displaystyle{\left(\coprod_{[n]\in\boldsymbol{\Delta}}X_n^{nd}\times\left|\Delta^n\right|\right)/\tildesim_f}} \\
	{[(\alpha_*(y),p)]} & {[(y,\alpha^*(p))]_f}
	\arrow["\varphi", from=1-1, to=1-2]
	\arrow["\pi"', from=1-1, to=2-1]
	\arrow["{\pi_{nd}}", from=1-2, to=2-2]
	\arrow[dashed, " \psi "',  from=2-1, to=2-2]
	\arrow["\psi"', maps to, from=3-1, to=3-2]
	\end{tikzcd}\]
	as the unique dashed arrow, rendering the diagram commutative. Since each $X_n^{nd} \times \left| \Delta^n \right|$ admits a continuous inclusion functor into $X_n\times \left| \Delta^n \right| $, by universality these assemble into the hooked solid arrow in the diagram
	\[\begin{tikzcd}
	{\displaystyle{\coprod_{[n]\in\boldsymbol{\Delta}}X_n\times\left|\Delta^n\right|}} & {\displaystyle{\coprod_{[n]\in\boldsymbol{\Delta}}X_n^{nd}\times\left|\Delta^n\right|}} \\
	{\displaystyle{\left(\coprod_{[n]\in\boldsymbol{\Delta}}X_n\times\left|\Delta^n\right|\right)/\tildesim}} & {\displaystyle{\left(\coprod_{[n]\in\boldsymbol{\Delta}}X_n^{nd}\times\left|\Delta^n\right|\right)/\tildesim_f}}
	\arrow["\pi"', from=1-1, to=2-1]
	\arrow["i", hook, from=1-2, to=1-1]
	\arrow["{\pi_{nd}}", from=1-2, to=2-2]
	\arrow[" \eta "', dashed, from=2-2, to=2-1]
	\end{tikzcd}\]
	Note that $\pi$ identifies all elements related by $\sim $, and thus we would like to show that $(x,p)\sim_f (y,q)$ implies $i(x,p)\sim i(y,q)$. This is indeed true, but we defer the proof to \cref{B.1.4}. Since $\pi$ identifies elements under $\sim$, the composite $\pi i$ identifies elements under $\sim _f$, and by universality of quotient, there is a unique dashed arrow $\zeta$ as indicated, rendering the diagram commutative. We now show $\psi$ and $\zeta$ are inverses. One direction is clear. Let $(x,p)$ be a point in the (not quotiented) disjoint union of nondegenerate simplices. Then $x$ is already nondegenerate, so $\alpha^*=1$, giving us $\varphi i(x,p)=\varphi(x,p)=(x,p)$. For the converse, let $x=\alpha^*(y)$ be the representation of \cref{B.1.1}, and note that
	\[
		i\varphi(x,p)=i(y,\alpha_*(p))=(y,\alpha_*(p))\sim(\alpha^*(y),p)=(x,p)
	.\]
	Thus, $i\varphi$ descends to $\zeta\psi=1$ along the quotient. We thus have that these maps are inverses, so we must now show they are continuous.

	We already know $\zeta$ is continuous since inclusions are continuous. To show $\psi$ is continuous, recall that since $X_n$ has the discrete topology, the product is isomorphic to the copower
	\[
		X_n\times \left| \Delta^n \right| \cong X_n\cdot \left| \Delta^n\right| \cong \coprod_{x\in X_n} \left\{ x \right\} \times \Delta^n
	.\]
	Let $x=\alpha^*(y)$ as given in \cref{B.1.1}. Note that $x\mapsto y$ is continuous since $X_n$ is discrete, and $\alpha_*$ is continuous by functoriality of $\left| \Delta^{\bullet}  \right| $, so by universality of coproduct, given the diagram of solid arrows
	\[\begin{tikzcd}
	{\{x\}\times\left|\Delta^n\right|} & {\{y\}\times\left|\Delta^n\right|} \\
	{\displaystyle{\coprod_{x\in X_n}\{x\}\times\left|\Delta^n\right|}} & {\displaystyle{\coprod_{x\in X_n^{nd}}\{x\}\times\left|\Delta^n\right|}}
	\arrow["{(x\mapsto y)\times\alpha_*}", " z "', from=1-1, to=1-2]
	\arrow[from=1-1, to=2-1]
	\arrow[from=1-2, to=2-2]
	\arrow["{\coprod z}"', dashed, from=2-1, to=2-2]
	\end{tikzcd}\]
	there exists a unique dashed arrow as indicated, rendering the diagram commutative. Note that we have chosen to denote $(x\mapsto y)\times \alpha_*$ by $z$ for simplicity. Since the diagram is in $\Top $, note that the dashed arrow is in particular continuous. Recalling our isomorphism above, once again universality of coproduct assembles the solid arrows
	\[\begin{tikzcd}
	{X_n\times\left|\Delta^n\right|} & {X_n^{nd}\times\left|\Delta^n\right|} \\
	{\displaystyle{\coprod_{[n]\in\boldsymbol{\Delta}}X_n\times\left|\Delta^n\right|}} & {\displaystyle{\coprod_{[n]\in\boldsymbol{\Delta}}X_n^{nd}\times\left|\Delta^n\right|}}
	\arrow["{\coprod z}", from=1-1, to=1-2]
	\arrow[from=1-1, to=2-1]
	\arrow[from=1-2, to=2-2]
	\arrow["{\coprod\coprod z}"', from=2-1, to=2-2]
	\end{tikzcd}\]
	into a unique dashed arrow as indicated, rendering the diagram commutative. In particular, the dashed arrow $\coprod \coprod z$ is continuous. Inspection of these universals shows $\coprod\coprod z=\varphi$, and thus $\psi$ is continuous, completing the proof.
\end{proof}

\begin{lemma}\label{B.1.4}
	With notation as in \cref{B.1.3}, if $(x,p)\sim_f (y,q)$, then $i(x,p)\sim i(y,q)$.
\end{lemma}
\begin{proof}
	Since $\sim _f$ is generated by the rule stated in \cref{B.1.2}, it suffices to show the statement for the generating rule. Let $x=\alpha^*(y)=\beta^*(z)$, where $\alpha^*(y)$ is as given in \cref{B.1.1}, and $\beta^*$ is a composite of faces. Then $i(y,\alpha_*(p))=(y,\alpha_*(p))$ and $i(z,\beta_*(p))=(z,\beta_*(p))$. We have
	\[
		(z,\beta_*(p))\sim (\beta^*(z),p)=(x,p)=(\alpha^*(y),p)\sim (y,\alpha_*(p))
	.\]
\end{proof}

\begin{remark}\label{B.1.5}
	One may be concerned that we assumed the simplex $x$ was degenerate throughout $\cref{B.1.3}$, and never considered the case where $x$ is nondegenerate. If $x$ is nondegenerate, then $\alpha^*=1$, which is clealy epi since it is an isomorphism, and thus the same proof holds. The same logic applies to all scenarios in which one might be concerned we did not consider a nondegenerate case.
\end{remark}

